\cleardoublepage
\phantomsection
\addcontentsline{toc}{chapter}{KATA PENGANTAR}
\begin{center}
{\FontJudulHalaman\bfseries KATA PENGANTAR\par}
\end{center}
\vspace{12pt}
Puji syukur penulis panjatkan kepada Allah SWT sehingga laporan kerja praktik di Divisi Maintenance PT XYZ dapat disusun. Laporan ini membahas integrasi perangkat pengukuran, HMI, dan dashboard sistem monitoring mesin cutting dan panel listrik.

Kegiatan ini memperluas pemahaman penulis mengenai instrumentasi, komunikasi data, perangkat lunak, serta tanggung jawab pemeliharaan di industri. Penulis menyampaikan terima kasih kepada:

\begin{enumerate}[leftmargin=*,itemsep=2pt]

\item Orang tua dan keluarga atas dukungan moral dan material selama pelaksanaan kerja praktik.

\item Bapak Ahmad Zainudin, S.ST., M.T., Ph.D., selaku Kepala Departemen Teknik Elektro Politeknik Elektronika Negeri Surabaya.

\item Bapak \NamaKetuaProdi, selaku Kepala Program Studi \JenjangProdi\ \NamaProdi\ Politeknik Elektronika Negeri Surabaya.

\item Ibu \NamaKoordinatorKP, selaku Koordinator Kerja Praktik Program Studi \JenjangProdi\ \NamaProdi\ Politeknik Elektronika Negeri Surabaya.

\item Bapak Dr. Ir. Agus Indra Gunawan, S.T., M.Sc., selaku dosen pembimbing, atas arahan dan masukan dalam pelaksanaan kegiatan serta penyusunan laporan kerja praktik.

\item PT XYZ atas kesempatan dan fasilitas yang diberikan selama pelaksanaan kerja praktik.

\item Bapak Dwi Jasmianto, selaku Senior Manager PT XYZ, atas dukungan terhadap pelaksanaan kegiatan kerja praktik.

\item Bapak Imsyar Lubis Tandika, selaku pembimbing lapangan, atas bimbingan teknis dan evaluasi terhadap hasil pekerjaan.

\item Personel Divisi Maintenance, staf, dan direksi PT XYZ atas pendampingan, kerja sama, serta dukungan selama kerja praktik.

\item Rekan-rekan yang memberikan bantuan dan dukungan selama pelaksanaan kegiatan serta penyusunan laporan.

\end{enumerate}

Penulis mengharapkan kritik dan saran untuk penyempurnaan laporan. Semoga laporan ini bermanfaat sebagai referensi penerapan sistem monitoring mesin.

\begin{flushright}Surabaya, \hspace{2.5cm}\\[12pt]Penulis\end{flushright}
